% MIR Medical Intelligence — IEEE-Style Research Paper Template
% Double-column manuscript skeleton for a specific study.
% Replace the placeholders and tailor the section structure to the target
% IEEE journal/conference author instructions before submission.

\documentclass[10pt,journal]{IEEEtran}

\usepackage[T1]{fontenc}
\usepackage[utf8]{inputenc}
\usepackage{lmodern}
\usepackage{microtype}
\usepackage{graphicx}
\usepackage{booktabs}
\usepackage{multirow}
\usepackage{amsmath,amssymb,amsfonts}
\usepackage{siunitx}
\usepackage{enumitem}
\usepackage{xcolor}
\usepackage{hyperref}
\usepackage[nameinlink,noabbrev]{cleveref}
\usepackage{url}
\usepackage{float}
\usepackage{placeins}
\usepackage{threeparttable}
\usepackage{tabularx}

\hypersetup{
    colorlinks=true,
    linkcolor=blue,
    citecolor=blue,
    urlcolor=blue,
    pdftitle={MIR Medical Intelligence — Research Article},
    pdfauthor={Mir AbdulRehman and Sababa Ateeq}
}

\setlist[itemize]{leftmargin=*}
\setlist[enumerate]{leftmargin=*}

% ---------- Paper metadata ----------
\newcommand{\PaperTitle}{Insert Study-Specific Title Here}
\newcommand{\MirName}{Mir AbdulRehman}
\newcommand{\SababaName}{Sababa Ateeq}
\newcommand{\MirRole}{Co-Research Lead --- AI/ML \& Software Engineering}
\newcommand{\SababaRole}{Co-Research Lead --- Biomedical \& Medical AI}
\newcommand{\MirORCID}{0000-0000-0000-0000}% Replace
\newcommand{\SababaORCID}{0000-0000-0000-0000}% Replace
\newcommand{\AffiliationOne}{MIR Medical Intelligence / Independent Research Program}
\newcommand{\AffiliationTwo}{Department or Institution, City, Country}% Replace where applicable
\newcommand{\CorrespondingEmail}{research@example.org}% Replace
\newcommand{\StudyKeywords}{medical AI, laboratory medicine, machine learning, clinical prediction, uncertainty, explainable AI}

\title{\PaperTitle}

\author{\MirName, \MirRole, ORCID: \href{https://orcid.org/\MirORCID}{\MirORCID},\\
\SababaName, \SababaRole, ORCID: \href{https://orcid.org/\SababaORCID}{\SababaORCID}\\
\AffiliationOne\\
\AffiliationTwo\\
Corresponding author: \href{mailto:\CorrespondingEmail}{\CorrespondingEmail}}

\begin{document}
\maketitle

\begin{abstract}
% Structured or narrative abstract as required by the target publication.
% State: Background/Objectives; Methods; Results; Conclusions; and, where useful,
% dataset, population, key metric, and principal scientific contribution.
\end{abstract}

\begin{IEEEkeywords}
\StudyKeywords
\end{IEEEkeywords}

\section{Introduction}
\subsection{Clinical/scientific background}
\subsection{Problem statement}
\subsection{Existing evidence}
\subsection{Research gap}
\subsection{Study objective}
\subsection{Hypothesis}
\subsection{Contributions}
% End the introduction with a concise statement of what this paper contributes.

\section{Related Work}
\subsection{Clinical prediction and medical AI}
\subsection{Laboratory-based modeling}
\subsection{Missing-data and missingness-aware methods}
\subsection{Uncertainty and calibration}
\subsection{Explainability}
\subsection{Adaptive evidence acquisition}
\subsection{Position of the present study}

\section{Materials and Methods}
\subsection{Study design}
\subsection{Research setting and data sources}
\subsection{Dataset description}
\subsection{Ethics and data governance}
\subsection{Population and cohort construction}
\subsection{Inclusion and exclusion criteria}
\subsection{Index date and observation window}
\subsection{Outcome definition}
\subsection{Predictor variables}
\subsection{Laboratory representation}
\subsection{Reference ranges and units}
\subsection{Temporal feature construction}
\subsection{Missingness representation}
\subsection{Data preprocessing}
\subsection{Data leakage prevention}
\subsection{Train/validation/test strategy}
\subsection{External validation}

\section{Model Development}
\subsection{Reference/rule-based baseline}
\subsection{Statistical baselines}
\subsection{Machine learning models}
\subsection{Deep learning models}
\subsection{Ensemble models}
\subsection{Hyperparameter optimization}
\subsection{Class imbalance handling}
\subsection{Calibration methods}
\subsection{Uncertainty quantification}
\subsection{Ablation studies}

\section{Adaptive Evidence Acquisition Methodology}
% Remove this whole section for papers that do not study adaptive testing.
\subsection{Problem formulation}
\subsection{Candidate evidence set}
\subsection{Evidence-value metric}
\subsection{Sequential decision policy}
\subsection{Stopping criterion}
\subsection{Clinician-in-the-loop constraint}
\subsection{Evaluation design}

\section{Evaluation and Statistical Analysis}
\subsection{Primary endpoint}
\subsection{Secondary endpoints}
\subsection{Discrimination metrics}
\subsection{Calibration metrics}
\subsection{Clinical utility}
\subsection{Survival metrics}
\subsection{Anomaly-detection metrics}
\subsection{Uncertainty metrics}
\subsection{Subgroup analysis}
\subsection{Sensitivity analyses}
\subsection{Statistical tests and confidence intervals}
\subsection{Model comparison procedure}
\subsection{Multiplicity and statistical safeguards}

\section{Explainability and Clinical Communication}
\subsection{Global feature importance}
\subsection{Local explanations}
\subsection{Attribution methods}
\subsection{Evidence coverage}
\subsection{Uncertainty communication}
\subsection{Inconclusive and abstention behavior}
\subsection{Clinical-language policy}

\section{Results}
\subsection{Cohort characteristics}
\subsection{Data completeness and missingness}
\subsection{Primary performance results}
\subsection{Calibration results}
\subsection{Clinical utility results}
\subsection{Subgroup results}
\subsection{External-validation results}
\subsection{Ablation results}
\subsection{Uncertainty results}
\subsection{Adaptive evidence-acquisition results}

\begin{table*}[t]
\centering
\caption{Primary model performance. Replace placeholder values with estimates and confidence intervals.}
\label{tab:main-results}
\begin{tabular}{lcccccc}
\toprule
Model & AUROC & AUPRC & Sensitivity & Specificity & Brier & Calibration slope \\
\midrule
Baseline & -- & -- & -- & -- & -- & -- \\
Logistic/Elastic Net & -- & -- & -- & -- & -- & -- \\
Tree model & -- & -- & -- & -- & -- & -- \\
Deep model & -- & -- & -- & -- & -- & -- \\
Proposed model & -- & -- & -- & -- & -- & -- \\
\bottomrule
\end{tabular}
\end{table*}

\begin{figure}[t]
    \centering
    \fbox{\rule{0pt}{1.7in}\rule{0.95\columnwidth}{0pt}}
    \caption{Placeholder for the principal model/analysis figure.}
    \label{fig:main}
\end{figure}

\section{Discussion}
\subsection{Principal findings}
\subsection{Comparison with prior work}
\subsection{Scientific interpretation}
\subsection{Clinical interpretation}
\subsection{Why the method succeeds or fails}
\subsection{Implications for missingness-aware intelligence}
\subsection{Implications for uncertainty-aware decision support}
\subsection{Implications for future adaptive evidence acquisition}

\section{Limitations}
\subsection{Dataset limitations}
\subsection{Selection and indication bias}
\subsection{Missingness limitations}
\subsection{Generalizability}
\subsection{Measurement and reference-range limitations}
\subsection{Causal-inference limitations}
\subsection{Clinical-deployment limitations}
\subsection{Residual model risk}

\section{Clinical Safety, Ethics, and Regulatory Considerations}
\subsection{Research-use boundary}
\subsection{Human oversight}
\subsection{Risk management}
\subsection{Data protection and privacy}
\subsection{Regulatory considerations}
\subsection{Potential failure modes}
\subsection{Safe-use recommendations}

\section{Reproducibility and Research Transparency}
\subsection{Code availability}
\subsection{Data availability}
\subsection{Dataset access restrictions}
\subsection{Software and environment}
\subsection{Versioning and archival release}
\subsection{Experiment registry}
\subsection{Model card / supplementary material}

\section{Conclusion}
% State the principal finding, scientific contribution, limits, and next research step.

\section*{Funding}
% State funding or: "This research received no external funding."

\section*{Conflicts of Interest}
% Declare conflicts or state none.

\section*{Ethics Statement}
% State ethics review / exemption / dataset governance as applicable.

\section*{Author Contributions}
\textbf{Mir AbdulRehman:} Conceptualization; Methodology; Software; Data Curation; Formal Analysis; Investigation; Validation; Visualization; Writing---Original Draft; Writing---Review \& Editing.\\
\textbf{Sababa Ateeq:} Conceptualization; Methodology; Biomedical Interpretation; Investigation; Validation; Writing---Original Draft; Writing---Review \& Editing.\\
% Replace with the actual contribution record for this paper. Use the CRediT taxonomy where appropriate.

\section*{Acknowledgment}
% Acknowledge contributors who do not meet the applicable authorship criteria.

\bibliographystyle{IEEEtran}
\bibliography{../shared/mir_references}

\appendices
\section{Supplementary Methods}
% Optional.
\subsection{Preprocessing details}
\subsection{Hyperparameter search}
\subsection{Additional ablations}

\section{Additional Results}
% Optional.

\end{document}
